\documentclass[11pt,a4paper]{article}
\usepackage[T1]{fontenc}
\usepackage{lmodern}
\usepackage[margin=1in]{geometry}
\usepackage{amsmath,amssymb,amsthm,mathtools,booktabs,array}
\usepackage{microtype,enumitem,xcolor,needspace}
\usepackage{hyperref}
\usepackage{xurl}
\hypersetup{colorlinks=true,linkcolor=blue!45!black,citecolor=blue!45!black,
 urlcolor=blue!45!black,
 pdftitle={A pure-state matched-basis information inequality violation in dimension five},
 pdfauthor={Dev D. Goyal},
 pdfsubject={Exact matched-basis information inequality and formal verification}}
\setlength{\parindent}{0pt}
\setlength{\parskip}{0.35em}
\newtheorem{theorem}{Theorem}
\title{A pure-state matched-basis information\\inequality violation in dimension five}
\author{Dev D. Goyal}
\date{28 October 2026}
\begin{document}
\maketitle
\paragraph{Autonomous AI research disclosure.}
This work is fully autonomous agentic research. Automated systems produced the research, proofs, formalization, and manuscript preparation. Dev D. Goyal is responsible for publication. No expert validation or peer review is claimed.

\begin{abstract}
This note defines a matched-identical-basis information comparison: the
minimum sum from five of six mutually unbiased local settings is compared
with the quantum mutual information of the state. In local dimension five,
the normalized pure state
$5^{-1/2}\sum_{x\in\mathbb F_5}|x\rangle|2x\rangle$ has quantum mutual
information $2\log5$. When both parties use the same computational basis
or any of the five specified quadratic bases, the joint measurement table
is a uniform permutation table with mutual information $\log5$.
Every selection of five settings therefore has sum $5\log5$, exceeding
the quantum-information budget by $3\log5$, with ratio $5/2$.
The proof verifies all six orthonormal bases, mutual unbiasedness, the
actual Born probabilities, purity, and both reduced states. The basis
family and identical-local-basis convention are explicit. Prior mixed-state
ECQC counterexamples receive separate credit. Correspondence with an
external ECQC proposal and historical novelty remain unverified.
\end{abstract}

\section{The question and prior work}
Classical correlations obtained from several complementary measurements
need not fit within a proposed quantum-information budget. The question
here is whether the explicitly defined matched-basis bound below holds
on pure states. The construction below shows a violation for a fixed complete
six-basis family in dimension five. The mechanism is the exact identity
$1+2^2=0$ in $\mathbb F_5$, which cancels every quadratic phase in the
same-basis measurement amplitude.

For a bipartite density operator \(\rho\), write

\[
 I(A:B)_\rho=S(\rho_A)+S(\rho_B)-S(\rho),
 \qquad S(\rho)=-\operatorname{Tr}(\rho\log\rho).
\]

For a local basis measurement \(M\) performed by both parties, let
\(I(M^A:M^B)\) denote the Shannon mutual information of its joint Born
distribution. All logarithms in this note are natural; division by
\(\log2\) gives bits. The choice of base has no effect on the
inequality.

For a finite probability distribution, $H(P)=-\sum_x P(x)\log P(x)$,
with $0\log0=0$. Classical mutual information is
$H(P^A)+H(P^B)-H(P)$. Two orthonormal bases are mutually unbiased when
every squared cross-basis overlap is $1/d$.

For a fixed complete collection \(\mathcal M\) of \(d+1\) mutually
unbiased bases, define the score and the inequality to be tested by

\[
 E_d(\rho;\mathcal M):=
 \min_{\substack{\mathcal S\subset\mathcal M\\|\mathcal S|=d}}
 \sum_{M\in\mathcal S}I(M^A:M^B)
 \ \le\ I(A:B)_\rho.
 \tag{1}
\]

The minimum discards one largest measurement mutual information. The
collection is fixed; this definition does not optimize over other basis
collections or local changes of basis.

The motivating ECQC proposal and pure-state question are attributed to
Iqbal~\cite{Iqbal}. Wang, Wang, and Chen receive credit for the earlier
mixed-state ECQC counterexamples attributed to their cited
preprint~\cite{WangECQC}. These source comparisons, the precise published
optimization, Bob's measurement convention, and current versions have
not been verified here. A refutation of an external statement follows
only if its definitions and quantifiers agree with (1) and the measurement
model below. No unrestricted ECQC priority claim is made.

The theorem below violates (1) using identical local basis vectors on
both parties. Bob uses the same basis vectors as Alice in each setting.
Replacing his basis by a separately conjugated family changes the model.

\section{Exact construction and theorem}

Identify indices with \(\mathbb F_5=\mathbb Z/5\mathbb Z\), and let

\[
 \omega=e^{2\pi i/5}.
\]

Define the normalized pure state

\[
 |\Psi\rangle=\frac1{\sqrt5}\sum_{x\in\mathbb F_5}|x\rangle_A|2x\rangle_B.
 \tag{2}
\]
The same vector in computational notation is
\[
 |\Psi\rangle=\frac{|00\rangle+|12\rangle+|24\rangle+|31\rangle+|43\rangle}{\sqrt5}.
\]

Take the computational basis \(M_\infty\) and the five quadratic bases

\[
 M_a=\{|a,k\rangle:k\in\mathbb F_5\},\qquad
 |a,k\rangle=\frac1{\sqrt5}\sum_{x\in\mathbb F_5}
 \omega^{a x^2+kx}|x\rangle,
 \quad a\in\mathbb F_5.
 \tag{3}
\]

The set \(\mathcal M=\{M_\infty,M_0,M_1,M_2,M_3,M_4\}\) has six
elements. Subscripts on the quadratic bases in this note are field
indices. Section 4 gives an equivalent matrix representation.

\begin{theorem}[Pure-state matched-basis violation]\label{thm:main}
The five quadratic bases in (3) and the computational basis form a complete
set of mutually unbiased orthonormal bases. For the state (2), every same-basis
measurement has joint distribution

\[
 P_M(i,j)=
 \begin{cases}
  1/5,&j=2i\pmod5,\\
  0,&\text{otherwise}.
 \end{cases}
 \tag{4}
\]

Consequently,

\[
 I(M^A:M^B)=\log5\quad(M\in\mathcal M),\qquad
 I(A:B)_{|\Psi\rangle\langle\Psi|}=2\log5,
 \tag{5}
\]

and

\[
 E_5(|\Psi\rangle\langle\Psi|;\mathcal M)
   =5\log5>2\log5.
 \tag{6}
\]

\end{theorem}

This is a counterexample to the displayed universal inequality
on pure states. The excess is exactly \(3\log5\), and the ratio of the
two sides is \(5/2\). No assertion of minimal counterexample dimension
is made.

\section{Proof of Theorem~\ref{thm:main}}

\subsection{A fifth-root identity}

For \(c\in\mathbb F_5\),

\[
 \sum_{x\in\mathbb F_5}\omega^{cx}=
 \begin{cases}5,&c=0,\\0,&c\ne0.\end{cases}
 \tag{7}
\]

The zero case is immediate. In the other case, multiplication by \(c\)
permutes \(\mathbb F_5\), and

\[
 1+\omega+\omega^2+\omega^3+\omega^4=0.
\]

All exponent calculations below occur modulo five.

\subsection{Orthonormality and mutual unbiasedness}

For a fixed \(a\),

\[
 \langle a,k|a,l\rangle
 =\frac15\sum_x\omega^{(l-k)x}
 =\delta_{kl}
\]

by (7). Thus each \(M_a\) is an orthonormal basis. Every computational
coordinate of every vector in \(M_a\) has squared magnitude \(1/5\),
giving mutual unbiasedness with \(M_\infty\).

For two distinct quadratic bases, set \(c=b-a\ne0\) and \(e=l-k\). Their
unnormalized overlap is

\[
 G(c,e)=\sum_x\omega^{cx^2+ex}.
\]

Writing \(x=y+u\), expanding, and then applying (7),

\[
\begin{aligned}
 |G(c,e)|^2
 &=\sum_{u,y}\omega^{c((y+u)^2-y^2)+eu}\\
 &=\sum_u\omega^{cu^2+eu}\sum_y\omega^{2cuy}\\
 &=5.
\end{aligned}
\]

Indeed, \(2c\ne0\) in \(\mathbb F_5\), so the inner sum is zero unless
\(u=0\), when it equals five. After the two factors \(1/\sqrt5\) in the
basis vectors, the squared overlap is \(5/25=1/5\). All six bases are
therefore pairwise mutually unbiased.

\subsection{Normalization and reduced states}

The map \(x\mapsto2x\) is a permutation of \(\mathbb F_5\), and the five
computational terms in (2) are orthogonal. Hence
\(\langle\Psi|\Psi\rangle=1\).

Let \(\rho=|\Psi\rangle\langle\Psi|\). Tracing out either subsystem
gives

\[
 \rho_A=\rho_B=\frac15I_5.
 \tag{8}
\]

\Needspace{5\baselineskip}
For example,

\[
 \operatorname{Tr}_B\rho
 =\frac15\sum_{x,y}|x\rangle\langle y|\,
   \langle2y|2x\rangle
 =\frac15\sum_x|x\rangle\langle x|.
\]

The other partial trace follows in the same way, using that
multiplication by two is bijective.

\subsection{Every joint measurement table}

In the computational basis, (4) follows immediately from (2).

For a quadratic basis \(M_a\), the joint amplitude is

\[
\begin{aligned}
 (\langle a,i|\otimes\langle a,j|)|\Psi\rangle
 &=\frac1{5\sqrt5}\sum_x
   \omega^{-a x^2-i x-a(2x)^2-j(2x)}\\
 &=\frac1{5\sqrt5}\sum_x\omega^{-(i+2j)x}.
\end{aligned}
 \tag{9}
\]

The second equality uses the exact field identity \(1+2^2=0\). By (7),
the amplitude is \(1/\sqrt5\) when \(i+2j=0\), and zero otherwise. Since
\(2^2=-1\), the equation \(i+2j=0\) is equivalent to \(j=2i\). Taking
the squared magnitude proves (4) for every quadratic setting.

With row and column order \(0,1,2,3,4\), the common table is explicitly

\[
 P_M=\frac15
 \begin{pmatrix}
 1&0&0&0&0\\
 0&0&1&0&0\\
 0&0&0&0&1\\
 0&1&0&0&0\\
 0&0&0&1&0
 \end{pmatrix}.
 \tag{10}
\]

\subsection{Entropies and the strict violation}

Both marginals of (10) are uniform. The joint table has exactly five
nonzero entries, all \(1/5\). Therefore

\[
 H(P_M^A)=H(P_M^B)=H(P_M)=\log5,
\]

and

\[
 I(M^A:M^B)=H(P_M^A)+H(P_M^B)-H(P_M)=\log5.
\]

The state \(\rho\) is a rank-one orthogonal projector, so its nonzero
spectrum is \((1)\) and \(S(\rho)=0\). Equation (8) gives
\(S(\rho_A)=S(\rho_B)=\log5\), proving the quantum value in (5).

Every five-element subset of the six measurement settings has the same
sum, \(5\log5\); thus the minimum in (1) is exactly that value. Since
\(\log5>0\), equation (6) is strict. This completes the proof.

In bits, the six individual scores are each approximately
\(2.3219280949\), the matched-basis score is approximately \(11.6096404744\), and
the quantum mutual information is approximately \(4.6438561898\). These
decimal values are illustrations of the exact formulas, not inputs to
the proof.

\Needspace{8\baselineskip}
\section{Equivalent matrix representation and source comparison}

The displayed bases can also be written using the Fourier matrix \(F\),
the identity, and powers of

\[
 D=\operatorname{diag}(1,\omega,\omega^4,\omega^4,\omega),
 \qquad F_{xk}=\omega^{xk}/\sqrt5.
\]

Because \(x^2\) modulo five is \(0,1,4,4,1\), the columns of \(D^aF\)
are exactly (3). The correspondence is:

\begin{center}
\begin{tabular}{@{}lll@{}}
\toprule\noalign{}
Comparison label & Matrix & Label in this note \\
\midrule\noalign{}
\(M_0\) & \(F\) & \(M_0\) \\
\(M_1\) & \(I_5\) & \(M_\infty\) \\
\(M_2\) & \(DF\) & \(M_1\) \\
\(M_3\) & \(D^2F\) & \(M_2\) \\
\(M_4\) & \(D^3F\) & \(M_3\) \\
\(M_5\) & \(D^4F\) & \(M_4\) \\
\bottomrule
\end{tabular}
\end{center}

Reindexing these six bases has no effect on the minimum over five-element
subsets. The comparison labels record the convention attributed to
Iqbal's Section 4.5.2~\cite{Iqbal}; their identity with the published
vectors and local measurement rule remains unverified. The matrix
identities in this section follow directly from (3).

The state differs from the conventional \(5^{-1/2}\sum_x|x,x\rangle\) by
a fixed local permutation. This distinction matters because the
measurement bases are fixed. The cancellation in (9) comes from the
permutation in the state. It does not change either party's measurement
rule.

\section{Scope and attribution}
The construction gives an exact pure-state violation for the displayed
matched measurement family. It does not assert the smallest possible
counterexample dimension. It also does not claim the first disproof of
unrestricted ECQC: the mixed classical-classical counterexamples of
Wang, Wang, and Chen are prior work~\cite{WangECQC}.

Original two-basis complementary quantum correlation (CQC) is a separate
question. Related contributions are attributed to Wang, Wang, and Chen
\cite{WangQubits} and Qian \cite{Qian}. These later-work summaries and
bibliographic details remain unverified. These attributions do not turn the present example into a new
resolution of original CQC. The result here is the exact
pure-state construction for the specified matched-basis settings. Its historical
priority and the comparison with current source versions remain provisional.

The state differs from the usual maximally entangled state by a local
permutation, while the measurement bases remain fixed. This distinction
is essential to the phase cancellation. The result makes no claim about
quantum capacity, cryptographic security, experimental implementation,
or a classification of all dimensions or measurement families.

\Needspace{10\baselineskip}
\section{Formal verification and reproducibility}
The formal project uses Lean 4.19.0 and mathlib revision
\begin{center}
\texttt{\small c44e0c8ee63ca166450922a373c7409c5d26b00b}.
\end{center}
The state is an actual normalized vector; its density matrix is its outer
product. Both reduced matrices are partial traces. Von Neumann entropy
is defined through continuous functional calculus of the matrices.
The measurement probabilities are squared magnitudes of the actual Born
amplitudes, then proved equal to the permutation table. Thus the entropy
calculation is connected to the quantum state and measurements, rather
than assumed for an unrelated table.

The principal declarations, in namespace \texttt{ECQC5}, are
\begin{verbatim}
pure_state_ECQC_counterexample
no_actual_five_setting_budget
actualScore_bits_violation
\end{verbatim}
The complete source was freshly rebuilt with Lean 4.19.0 on 28 October
2026. A separate audit with \texttt{-{}-trust=0} checked eight named
endpoints, including basis unitarity, mutual unbiasedness, the Born-table
identity, quantum mutual information, and density-matrix positivity.
The build and audit succeeded. All eight reports contain only
\texttt{propext}, \texttt{Classical.choice}, and \texttt{Quot.sound}.
The proof sources and pins were unchanged during checking. A declaration
map and reproduction commands accompany this paper. Formal checking does
not establish historical novelty, source attribution, expert validation,
or an external conjecture identification.

A supplementary exact cyclotomic arithmetic check covered 125 orthogonality
identities, 375 unbiasedness identities, 125 quadratic Born probabilities,
and 100 negative-control checks for the ordinary maximally entangled
state. These are finite exact checks. They support the explicit arithmetic
and do not replace the formal density-matrix and entropy proof.
The all-prime companion note proves an unbounded-ratio extension by a
conventional argument. The Lean package covers dimension five only.

\bibliographystyle{plain}
\bibliography{references}
\end{document}
